\newpage
\section[Modelling Exercise: Pyrite oxidation during deep well injection]{Modelling Exercise: Pyrite oxidation during deep well injection}

Managed aquifer recharge is increasingly used to enhance the sustainable development of water supplies. Common recharge techniques include aquifer storage and recovery (ASR), infiltration ponds, river bank filtration and deep-well injection. Following recharge the water quality of the injectant is typically altered by a multitude of geochemical processes during subsurface passage and storage. Relevant geochemical processes that affect the major ion chemistry include microbially mediated redox reactions, mineral dissolution/precipitation, sorption and ion-exchange. The hydrochemical conditions and changes that occur under these circumstances, in particular the temporal and spatial changes of pH and redox conditions, are in many cases the controlling factor for the fate of micropollutants such as herbicides and pharmaceuticals. Similarly, changes in mineralogical composition such as dissolution and precipitation of iron- or aluminiumoxides may affect the mobility of trace metals as well as the attachment and subsequent decay of pathogenic viruses. Laboratory and field-scale experimental studies are aimed at investigating such processes under controlled conditions and to eventually develop a better qualitative and quantitative understanding of their complex interactions, both site-specific and at a fundamental level. \
\cite{prommer_stuyfzand_2005} carried out a reactive transport 
modelling study to analyse the data collected during a deep well 
injection experiment in an anaerobic, pyritic aquifer near Someren in Southern Netherlands. 

This exercise replicates some of the key processes that were identified
to influence water quality changes during subsurface passage. 
Pyrite oxidation will be defined as a kinetic process in which the
reaction rate depends on the water temperature. 
For simplicity we include temperature as a seperate aqueous (mobile) component
called {\color{red} \textbf{Tmp}}.
The reaction rate expression expression has been programmed such 
that the value of the component {\color{red} \textbf{Tmp}} is read
and used during the computation of the reaction rate. 

\subsection[Setting up the flow problem]{Setting up the flow problem}

To simplify the original, fully three-dimensional model for this exercise, 
the flow model is defined as a two-dimensional model for a single stratigraphic layer. Aerobic surface water is injected through an injection well and extracted 
at an extraction well located 100 away from the injection well. 

\begin{itemize}

\item Use the data supplied in Table \ref{ex_model_grid} to set up the flow model. 
Note that the geometry is defined for a half-model due to the symmetry of the problem.
Neglect background groundwater flow and use for simplicity fixed head cells at the left 
(x = 0 m) and the right (x = 200 m) model boundary to allow 
exchange of water across these boundaries.

\item Check your simulated heads in the injection well (24.22 m) and in the recovery well (15.78 m).
You can also visualise the heads along a profile that goes across the injection and extraction well
by going to
{\color{blue}Observation} $\rightarrow$ {\color{blue}OBS} $\rightarrow$ {\color{blue}obs.1 Observation} 
and defining a line. 

\end{itemize}

\begin{table}
\caption{Flow and transport parameters used for the deep well injection problem.}
\begin{center}
\begin{tabular}{lc} \hline
Flow simulation                     & steady state \\
Total simulation time $(days)$       & 360  \\
Time steps length $(days)$       & 2 \\
Model \  length\ (column direction) \ $ (m)      $         &  200            \\
Model \  width\  (row direction) \ $ (m)      $         &   80            \\
Grid \ spacing\ $\Delta x $ (m)                &   10            \\
Grid \ spacing\ $\Delta y $ (m)                &   10            \\
Total Porosity\     $            $         &  0.35           \\
Effective Porosity\     $            $         &  0.35           \\
Horizontal hydraulic conductivity\ $(m)           $         &  10         \\  
Piezometric head upstream (left) boundary\ $(m)          $         &  20           \\ 
Piezometric head downstream (right) boundary\ $(m)        $         &  20           \\ 
Depth target aquifer $(m \ below \ sealevel)   $         &  between -300 and -310          \\ 
x-Position injection well $(m)   $         & 145          \\
y-Position injection well $(m)   $         &   5         \\
x-Position recovery well $(m)    $         &  55          \\
y-Position recovery well $(m)    $         &   5         \\
1/2 Flow rate injection well $(m^3/day)   $         &  500         \\
1/2 Flow rate recovery well $(m^3/day)    $         &  500         \\
Longitudinal dispersivity\ $(m)         $         &  1          \\ 
Transversal  dispersivity\ $(m)         $         &  0.1          \\ \hline

\end{tabular}
\label{ex_model_grid}
\end{center}
\end{table}

\section[Setting up the nonreactive transport problem]{Setting up the nonreactive transport problem}

\begin{itemize}

\item If the simulated heads agree, use MT3DMS to simulate a simple tracer experiment
to estimate the travel times between the injection and the recovery well. 
Set the tracer concentration in the background water to 0 and set the tracer concentration
in the injection well to 1. Include some observation points on the axis between injection and recovery well to allow the visualisation
of breakthrough curves.

\item What is your simulated travel time to the mid-point between 
injection and recovery well ? It should be approximately 20 days.
If your result differs, try to debug the problem before proceeding.

\item When does the recovery well receive 100\% of the injection concentration ? 

\end{itemize} 

\subsection[Setting up the reactive transport problem]{Setting up the reactive transport problem}

Once the tests with the nonreactive model indicate that the model is working properly,
proceed with the reactive transport model. 

\begin{itemize}

\item Set appropriate boundary conditions for the fluxes/concentrations across the model boundary. 

\item Copy the reaction database that was developed for this problem into the folder where the iPHT3D model is located. 

\item Activate the relevant aqueous components that are needed to simulate the reactive transport of the chemicals listed in Table 1 of \cite{prommer_stuyfzand_2005}. Also include  {\color{red} \textbf{Tmp}} (Temperature). 
For simplicity do not include  {\color{red} \textbf{DOC}} and ion exchange species. 
From the minerals listed add only  {\color{red} \textbf{Pyrite (kinetic version)}} 
and \\ {\color{red} \textbf{Fe(OH)3(a)}} to the reaction network.  

\item Define the initial concentrations (background water composition and mineral concentrations)
in the appropriate iPHT3D Tab. Use the 
water composition listed in Table 1 of \cite{prommer_stuyfzand_2005}.
For pyrite use an initial concentration of 0.05 mol/l bulk volume and define 
it under the {\color{blue} \textbf{Phase}} Tab.
Note that the oxidation rate of pyrite is, among other factors, depending on the 
pyrite concentration. Where pyrite is present at higher concentrations the reaction
will proceed faster. For the temperature {\color{red} \textbf{Tmp}} enter a value of 0.017, which
corresponds to 1/1000 of the temperature in Celcius.

\item Define the water composition of the injected water as $Solution 1$.
Use the water composition measured at the 21 January 1997, as
listed in Table 1 of \cite{prommer_stuyfzand_2005}. For the temperature add a
value of 0.0019 (= 1.9 C). For pe, which is not given in Table 1, use a value of 13.5.
 
\item Attribute this water composition to the injection well 
under {\color{blue}Spatial Attributes} $\rightarrow$ {\color{blue}PHT3D} $\rightarrow$ {\color{blue}PH} $\rightarrow$ {\color{blue}ph.4 Source \/ Sink} and define a zone 
(a single point in this instance) at the well location. 

\item Make sure that the reaction rate parameters 1 - 5 for the kinetically
controlled pyrite oxidation are set to values of 16, .67, 0.5, -.11 and 115, respectively.

\item Select the TVD scheme as advection package and define the dispersivities.

\item Run the reactive transport model.

\item Inspect the results: How far did oxygen penetrate into the aquifer at the end of the simulation time ?

\end{itemize} 


\section[Mobilization of arsenic]{Mobilization of arsenic}

In the next part or the exercise we are going to model the release of arsenic 
during oxidation of arsenopyrite and its subsequent sorption to ferrihydrite (HfO). 
In a first step we will model the dissolution of \emph{Arsenopyrite}. 
The rate expression that is used for \emph{Arsenopyrite} will simply be linked
to the rate of the computed pyrite oxidation:

{\color{red}
\begin{quote}
\#\#\#\#\#\#\#\#\#\#\#\#\#\#\#\#\#\#\#\#\#\#\#\#\#\#\\
Arsenopyrite \\
\#\#\#\#\#\#\#\#\#\#\#\#\#\#\#\#\#\#\#\#\#\#\#\#\#\#\\
-start \\
  6 moles = parm(1) * get(1) \\
200 save moles \\
-end \\
\end{quote}}

The stoichiometric ratio at which arsenopyrite will be dissolved, compared to pyrite, is determined
by {\color{red} \textbf{parm(1)}}. For example, if {\color{red} \textbf{parm(1)}} is set to
\emph{0.001} the dissolution of 1 mmol pyrite will be accompanied by 
the dissolution of 1 $\micro$mol \emph{Arsenopyrite}.   
This is achieved by including {\color{red} \textbf{put(1)}}
in the rate expression for pyrite:

{\color{red}
\begin{quote}
\#\#\#\#\#\#\#\#\#\#\#\#\#\#\#\#\#\#\#\#\#\#\#\#\#\# \\
Pyrite \\
\#\#\#\#\#\#\#\#\#\#\#\#\#\#\#\#\#\#\#\#\#\#\#\#\#\# \\
-start \\
... \\
... \\
... \\
90 moles = moles\_oxy + moles\_nitr + moles\_0 \\
100 if (moles $\>$ m) then moles = m \\
150 put(moles,1) \\
200 save moles \\
-end \\
\end{quote}}

We do now expand the last simulation and include additionally the two aqueous components 
\emph{As(5)} and \emph{As(3)} as well as the kinetic minerals \emph{Fe(OH)3(a)} and \emph{Arsenopyrite} in the reaction network. 

\begin{itemize}

\item Set the initial concentration for \emph{Arsenopyrite} to 0.01 mol/l$_{bulk}$.

\item Set the reaction rate constants for \emph{Fe(OH)3(a)} to 2 $\times$ 10$^{-13}$
and define that the ratio of \emph{Arsenopyrite} / \emph{Pyrite} dissolution 
is 0.001 (i.e., parm(1) = 0.001 for \emph{Arsenopyrite}).

\item Run the simulation and plot a breakthrough curve for \emph{As(5)} and \emph{As(3)} at the extraction well and at a location half-way between injection and the extraction well.
Save the breakthrough curves from the observation wells to ASCII files.

\item Now include surface complexation reactions by 
activating the two surface species \emph{Hfo\_w} and \emph{Hfo\_s}.
Then define the surface site densities for both of these species as 0.2 mol per mol of ferrihydrite 
and 0.005 mol per mol of ferrihydrite, respectively. Define those concentrations in the {\color{blue} \textbf{Site\_back}} column in the {\color{blue} \textbf{Surface}} Tab. 
Then add \emph{Fe(OH)3(a)} under the {\color{blue} \textbf{name}} 
column and {\color{blue} \textbf{kinetic}} under the {\color{blue} \textbf{switch}} column, 
both also in the {\color{blue} \textbf{Surface}} Tab.  
This will establish that the number of sorption sites is directly 
linked to the concentration of the (kinetically reacting) 
mineral phase \emph{Fe(OH)3(a)}. 
Finally, define among the global PHT3D parameters that the {\color{blue} \textbf{no\_edl}} 
option is used for the simulation.
To do so set in 
{\color{blue}PHT3D} $\rightarrow$ {\color{blue}PH} $\rightarrow$ {\color{blue}ph.6 specific options}
set {\color{red}SU\_OPT} to {\color{blue} \textbf{no\_edl}}. 
This means that electric double layer calculations in the surface complexation 
model will be omitted.  
    

\item Run the model and inspect again the \emph{As(5)} and \emph{As(3)} breakthrough curves at the extraction well

\item Compare the new results with the previous results (model run without surface complexation).

\end{itemize}




\section[Seasonally changing redox zonation]{Seasonally changing redox zonation}

In the following part of the exercise we attempt to mimic the effects of a seasonally changing
injection water temperature. This will result in a dynamically changing redox zonation.
To simulate this effect proceed as follows:

\begin{itemize}

\item To effectively consider the transiently changing water compositions, iPHT3D allows to 
efficiently create multiple water compositions.
to create the 12 different solutions that only differ with respect to their temperature allow first allow the model to have 12 solutions. 
To do so set in 
{\color{blue}PHT3D} $\rightarrow$ {\color{blue}PH} $\rightarrow$ {\color{blue}ph.6 specific options}
set {\color{red}NB\_solu} to 12. After that re-import the database and open the chemistry dialog
where you can now define up to 12 solutions. Then copy solution 1 to the next 11 
(using ctrl+C and ctrl+v) and change the 
temperature to 8, 5, 9, 12, 15, 18, 22, 25, 21, 14, and 12, respectively. 
\\
Note that the heat transport approximation in this exercise is not very accurate
as it ignores the retardation that occurs as a result of heat transfer between
the injectant and the sediment matrix. 
While neglected here, it is possible to use MT3DMS/PHT3D to describe 
the relevant heat transport equation exactly. 
% Ref to Rui and Peter Bayer's papers

\end{itemize}


